\PassOptionsToPackage{mathlines}{lineno}
\documentclass[twocolumn]{aastex631}

\usepackage{amsmath}
\usepackage{amssymb}
\usepackage{graphicx}

\providecommand{\shorttitle}[1]{}
\providecommand{\shortauthors}[1]{}
\shorttitle{Coma Cluster Dark Matter Profile with DESI}
\shortauthors{Jubair Hossain}

\begin{document}

\title{Dynamical Mass and Dark Matter Halo Modeling of the Coma Cluster with DESI Spectroscopy}

\correspondingauthor{Jubair Hossain}
\email{enigmatic.jihad@gmail.com}

\author{Jubair Hossain}
\affiliation{Independent Researcher}

\begin{abstract}
Measuring galaxy kinematics in rich clusters provides a direct way to map their gravitational potential and test cosmological predictions for dark matter halos. In this work, I use spectroscopic observations from the Dark Energy Spectroscopic Instrument (DESI) Early Data Release (Iron) to study the dynamics of the Coma Cluster (Abell 1656). Selecting confirmed galaxies within an angular radius of $\theta \le 2.0^\circ$ and applying an iterative $3\sigma$ velocity clipping yields a clean sample of 2,098 member galaxies. I divide the sample into eight radial bins of width $0.25^\circ$ and subtract spectrograph measurement errors in quadrature to construct an intrinsic line-of-sight velocity dispersion profile. The dispersion decreases from $1075.94 \pm 40.72\text{ km s}^{-1}$ in the core to $807.01 \pm 85.07\text{ km s}^{-1}$ near the virial boundary, with a secondary bump around $\theta \sim 1.4^\circ$ likely caused by the infalling NGC~4839 subgroup. I fit three halo models to the observed dispersion profile: a Singular Isothermal Sphere (SIS), a Soft-Core profile, and a Navarro-Frenk-White (NFW) profile. Model comparison using reduced chi-squared and the Bayesian Information Criterion (BIC) shows that the cuspy NFW profile provides the best fit ($\chi_\nu^2 = 0.69$, $\text{BIC} = 8.27$), while the flat SIS model is ruled out ($\chi_\nu^2 = 3.23$, $\Delta\text{BIC} = 16.42$). Integrating the NFW potential at an assumed distance of $D = 98.9\text{ Mpc}$ gives a scale radius of $r_s \approx 1.81\text{ Mpc}$ and a virial mass of $M_{\text{vir}} = (5.11 \pm 0.45) \times 10^{14}\,M_\odot$ at $r = 2.0\text{ Mpc}$, in good agreement with independent X-ray and weak lensing measurements. These findings support the view that cluster-scale dark matter halos follow the cuspy density profiles expected in standard $\Lambda\text{CDM}$ cosmology.
\end{abstract}

\keywords{Galaxy clusters (584) --- Dark matter (353) --- Kinematics (918) --- Velocity dispersion (1746) --- Coma Cluster (271)}

\section{Introduction} \label{sec:intro}

\subsection{Historical Background and Dark Matter in Coma}
Galaxy clusters are the largest gravitationally bound systems in the universe. Because their overall dynamics are dominated by gravity, measuring the motions of their member galaxies allows us to trace their underlying mass distributions. The idea of dark matter originally emerged from observations of the Coma Cluster (Abell 1656). In 1933, Fritz Zwicky measured the radial velocities of galaxies in Coma and found that they had a surprisingly high velocity dispersion, on the order of $\sigma_v \sim 1000\text{ km s}^{-1}$ \citep{zwicky1933}. 

Using the virial theorem, Zwicky estimated the total mass needed to keep the cluster together and found that it was many times larger than the mass accounted for by visible stars. He proposed that the cluster must contain a large amount of non-luminous matter (``dunkle Materie''). Subsequent multi-wavelength observations have confirmed this conclusion, showing that visible baryons in stars and hot gas make up only about $10\text{--}15\%$ of the total cluster mass, with the rest residing in dark matter.

\subsection{Dark Matter Halo Profiles: Cusps and Cores}
In the standard cosmological framework of Cold Dark Matter with a cosmological constant ($\Lambda\text{CDM}$), structure forms through the hierarchical collapse of collisionless dark matter. Cosmological $N$-body simulations predict that virialized dark matter halos have a universal cuspy density distribution, commonly described by the Navarro-Frenk-White (NFW) profile \citep{navarro1997}:
\begin{equation}
	\rho_{\text{NFW}}(r) = \frac{\rho_0}{\left(\frac{r}{r_s}\right) \left(1 + \frac{r}{r_s}\right)^2}
\end{equation}
where $r_s$ is a characteristic scale radius and $\rho_0$ is a density normalization. This profile has a steep inner slope ($\rho \propto r^{-1}$) and falls off as $\rho \propto r^{-3}$ at large distances.

In contrast, observations of low-surface-brightness dwarf galaxies often show flatter, cored central density distributions, creating the long-standing ``cusp-core'' debate. To test whether real halos are cuspy or cored, it is useful to test alternative models:
\begin{enumerate}
	\item \textbf{Singular Isothermal Sphere (SIS):} A simple model where density decays as $\rho \propto r^{-2}$, predicting a constant velocity dispersion at all radii ($\sigma_v(r) = \sigma_0$).
	\item \textbf{Soft-Core Profile:} A model where the central density flattens to a constant value within a core radius, preventing any central density singularity.
\end{enumerate}
Because galaxy clusters are massive, relatively young, and less affected by small-scale stellar feedback than dwarf galaxies, they provide an excellent setting to test whether dark matter halos naturally develop NFW-like cusps.

\subsection{Observational Opportunities with DESI}
Testing these halo models dynamically requires measuring the line-of-sight velocity dispersion of cluster galaxies across a wide radial range. Earlier dynamical studies of Coma \citep[e.g.,][]{kent1982, lokas2003} had to rely on samples of roughly 300 to 450 galaxies, which introduced noticeable statistical uncertainties in individual radial bins.

The Dark Energy Spectroscopic Instrument (DESI) provides an opportunity to revisit this problem. Mounted on the Mayall 4-meter telescope at Kitt Peak National Observatory, DESI utilizes 5,000 robotically positioned optical fibers to obtain dense spectroscopic coverage over large sky areas \citep{desi2023}. In this study, I use galaxy redshifts from the DESI Early Data Release (Iron) to isolate over 2,000 confirmed Coma members, extract an intrinsic velocity dispersion profile, and evaluate NFW, SIS, and soft-core halo models.

\section{Data Reduction and Sample Selection} \label{sec:sample_selection}

\subsection{Catalog Selection Flags}
Galaxy targets were extracted from the DESI Iron data release catalog (\texttt{zall-pix-iron.fits}). To ensure high purity, three standard selection flags were applied:
\begin{itemize}
	\item \texttt{SPECTYPE == 'GALAXY'}: Selects confirmed galaxy spectra, removing stars and quasars.
	\item \texttt{ZWARN == 0}: Ensures reliable spectroscopic redshift fits without pipeline warning flags.
	\item \texttt{ZCAT\_PRIMARY == True}: Selects the primary observation for targets observed multiple times, avoiding duplicate measurements.
\end{itemize}

\subsection{Velocity Calculation and Spatial Cone}
Because galaxies in massive clusters reach relative velocities of several thousand $\text{km s}^{-1}$, spectroscopic redshifts ($z$) were converted into line-of-sight Doppler velocities ($v_{\text{LOS}}$) using the relativistic Doppler formula:
\begin{equation}
	v_{\text{LOS}} = c \cdot \frac{(1+z)^2 - 1}{(1+z)^2 + 1}
\end{equation}
where $c$ is the speed of light. The measurement error in velocity was propagated as $v_{\text{err}} = c \cdot \delta z$.

A spatial cone search was centered on the central region of the Coma Cluster at Right Ascension $\alpha_0 = 194.953^\circ$ and Declination $\delta_0 = +27.981^\circ$. The selection radius was set to $\theta \le 2.0^\circ$ (which corresponds to roughly $3.5\text{ Mpc}$ at the distance of Coma), spanning a redshift range of $z \in [0.013, 0.033]$, corresponding to $v_{\text{LOS}} \approx 3{,}900\text{ to }9{,}900\text{ km s}^{-1}$.

\subsection{Interloper Rejection}
To remove foreground and background galaxies that fall along the line of sight but are not gravitationally bound to Coma, an iterative $3\sigma$ velocity clipping was applied:
\begin{equation}
	|v_{\text{LOS}} - v_{\text{median}}| \le 3\sigma_v
\end{equation}
The initial candidate sample had a median velocity of $v_{\text{median}} = 6948.18\text{ km s}^{-1}$ and a velocity dispersion of $\sigma_v = 960.54\text{ km s}^{-1}$. All 2,098 selected galaxies fell within the $3\sigma$ clipping range ($4104.28\text{ to }9716.75\text{ km s}^{-1}$). This confirms a final sample of $N = 2{,}098$ member galaxies within a $2^\circ$ radius.

\section{Velocity Dispersion Profile} \label{sec:velocity_dispersion}

\subsection{Radial Binning and Error Correction}
To construct the velocity dispersion profile, the 2,098 member galaxies were divided into eight concentric radial annuli of width $\Delta\theta = 0.25^\circ$ out to $\theta = 2.0^\circ$. For each annulus $k$, the radial distance was taken as the median angular separation of the galaxies in that bin, $\theta_k = \text{median}(\theta_i)$.

The observed velocity dispersion ($\sigma_{\text{obs}, k}$) was computed as the sample standard deviation with an $N_k - 1$ denominator. To determine the intrinsic physical velocity dispersion ($\sigma_{\text{int}, k}$), the mean squared spectrograph velocity uncertainties were subtracted in quadrature:
\begin{equation}
	\sigma_{\text{int}, k} = \sqrt{\max\left(0, \, \sigma_{\text{obs}, k}^2 - \langle v_{\text{err}}^2 \rangle_k\right)}
\end{equation}
The statistical uncertainty in each bin was computed using the standard error on the dispersion for a normal distribution:
\begin{equation}
	\delta\sigma_k = \frac{\sigma_{\text{int}, k}}{\sqrt{2(N_k - 1)}}
\end{equation}

\subsection{Profile Properties}
Table~\ref{tab:kinematic_profile} lists the number of galaxies, median radius, observed dispersion, intrinsic dispersion, and statistical error for each radial bin.

\begin{table*}[ht!]
	\centering
	\caption{Radial Velocity Dispersion Profile of the Coma Cluster}
	\label{tab:kinematic_profile}
	\begin{tabular}{cccccc}
		\hline \hline
		Bin Range & $N_{\text{gal}}$ & Median Radius $\theta_k$ & $\sigma_{\text{obs}}$ & Intrinsic $\sigma_{\text{int}}$ & Error $\delta\sigma$ \\
		          &                  & [deg]                    & [$\text{km s}^{-1}$]  & [$\text{km s}^{-1}$]            & [$\text{km s}^{-1}$] \\
		\hline
		$0.00^\circ - 0.25^\circ$ & 350 & $0.145^\circ$ & 1076.03 & 1075.94 & $\pm 40.72$ \\
		$0.25^\circ - 0.50^\circ$ & 360 & $0.374^\circ$ & 1003.21 & 1003.10 & $\pm 37.44$ \\
		$0.50^\circ - 0.75^\circ$ & 354 & $0.625^\circ$ & 965.97  & 965.85  & $\pm 36.35$ \\
		$0.75^\circ - 1.00^\circ$ & 345 & $0.875^\circ$ & 925.71  & 925.59  & $\pm 35.29$ \\
		$1.00^\circ - 1.25^\circ$ & 311 & $1.118^\circ$ & 860.76  & 860.68  & $\pm 34.57$ \\
		$1.25^\circ - 1.50^\circ$ & 210 & $1.369^\circ$ & 910.37  & 910.27  & $\pm 44.52$ \\
		$1.50^\circ - 1.75^\circ$ & 122 & $1.597^\circ$ & 930.02  & 929.95  & $\pm 59.78$ \\
		$1.75^\circ - 2.00^\circ$ & 46  & $1.845^\circ$ & 807.10  & 807.01  & $\pm 85.07$ \\
		\hline
	\end{tabular}
\end{table*}

The profile shows three clear features:
\begin{enumerate}
	\item \textbf{High central peak:} The dispersion reaches its highest value in the innermost bin ($\sigma_{\text{int}} = 1075.94 \pm 40.72\text{ km s}^{-1}$), corresponding to the deep central potential near the giant elliptical galaxies NGC~4874 and NGC~4889.
	\item \textbf{Steady decline:} Between $0.25^\circ$ and $1.25^\circ$, the dispersion decreases steadily from $1003.10\text{ km s}^{-1}$ to $860.68\text{ km s}^{-1}$.
	\item \textbf{Substructure around $1.4^\circ$:} A secondary increase is visible between $1.3^\circ$ and $1.6^\circ$ ($910\text{--}930\text{ km s}^{-1}$). This localized kinematic feature matches the location of the infalling NGC~4839 subgroup to the southwest of the cluster core. Beyond this feature, the dispersion drops to its lowest value of $807.01 \pm 85.07\text{ km s}^{-1}$ at the outer boundary.
\end{enumerate}

\section{Dark Matter Halo Modeling and Model Selection} \label{sec:models}

\subsection{Theoretical Dispersion Models}
To evaluate whether Coma is better described by an isothermal, cored, or cuspy dark matter halo, I fit the intrinsic velocity dispersion profile to three standard theoretical models:

\begin{enumerate}
	\item \textbf{Singular Isothermal Sphere (SIS):} A constant velocity dispersion across all radii:
	\begin{equation}
		\sigma_{\text{SIS}}(\theta) = \sigma_0
	\end{equation}

	\item \textbf{Soft-Core Profile:} A cored profile with a characteristic core radius $\theta_c$ that prevents a central density cusp:
	\begin{equation}
		\sigma_{\text{Core}}(\theta) = \frac{\sigma_{0,c}}{\left[1 + (\theta/\theta_c)^2\right]^{0.25}}
	\end{equation}

	\item \textbf{Navarro-Frenk-White (NFW) Profile:} The standard $\Lambda\text{CDM}$ cuspy model, characterized by a characteristic dispersion $\sigma_{\text{max}}$ and a scale angle $\theta_s$:
	\begin{equation}
		\sigma_{\text{NFW}}(\theta) = \sigma_{\text{max}} \cdot \sqrt{\frac{\theta_s}{\theta} \ln\left(1 + \frac{\theta}{\theta_s}\right)}
	\end{equation}
\end{enumerate}

\subsection{Fitting Method and Selection Metrics}
The parameters of each model were determined through weighted non-linear least-squares fitting using the 8 radial data points. To compare the models objectively, three standard statistical metrics were calculated:

\begin{enumerate}
	\item \textbf{Total and Reduced Chi-Squared ($\chi^2, \chi_\nu^2$):}
	\begin{equation}
		\chi^2 = \sum_{k=1}^N \left(\frac{\sigma_{\text{int}, k} - \sigma_{\text{model}}(\theta_k)}{\delta\sigma_k}\right)^2, \quad \chi_\nu^2 = \frac{\chi^2}{N - k_{\text{param}}}
	\end{equation}
	where $N = 8$ is the number of bins and $k_{\text{param}}$ is the number of free parameters.

	\item \textbf{Bayesian Information Criterion (BIC):}
	\begin{equation}
		\text{BIC} = \chi^2 + k_{\text{param}} \ln(N)
	\end{equation}
	The BIC balances the quality of fit against model complexity, penalizing models that introduce unnecessary parameters.

	\item \textbf{Relative Difference ($\Delta\text{BIC}$):}
	\begin{equation}
		\Delta\text{BIC} = \text{BIC}_{\text{model}} - \text{BIC}_{\text{NFW}}
	\end{equation}
	Following standard interpretations, a difference of $\Delta\text{BIC} > 2$ indicates positive evidence against a model, while $\Delta\text{BIC} > 10$ is considered very strong evidence against it.
\end{enumerate}

\subsection{Fitting Results}
The best-fit parameters, chi-squared values, and BIC results are summarized in Table~\ref{tab:model_comparison}.

\begin{table*}[ht!]
	\centering
	\caption{Comparison of Dark Matter Halo Models}
	\label{tab:model_comparison}
	\begin{tabular}{cccccccc}
		\hline \hline
		Model & Parameters & Best-Fit Value & $\chi^2$ & d.o.f. & $\chi_\nu^2$ & BIC & $\Delta\text{BIC}$ \\
		\hline
		SIS       & $\sigma_0$                     & $946.99 \pm 14.15\text{ km s}^{-1}$                               & 22.61 & 7 & 3.23 & 24.69 & $+16.42$ \\
		Soft-Core & $\sigma_{0,c}, \theta_c$       & $1027.35\text{ km s}^{-1}, 1.414^\circ$                           & 7.47  & 6 & 1.24 & 11.63 & $+3.36$  \\
		NFW       & $\sigma_{\text{max}}, \theta_s$ & $1093.30 \pm 42.66\text{ km s}^{-1}, 1.047^\circ \pm 0.353^\circ$ & 4.11  & 6 & 0.69 & 8.27  & $0.00$   \\
		\hline
	\end{tabular}
\end{table*}

\begin{figure}[ht!]
	\centering
	\plotone{"Model Selection (SIS vs Cored vs NFW).png"}
	\caption{Measured velocity dispersion profile of the Coma Cluster from DESI data (black points with error bars) compared to the three fitted models: SIS (blue dashed line), Soft-Core (orange dash-dotted line), and NFW (green solid line). The NFW model provides the best overall description of the data ($\chi_\nu^2 = 0.69$).}
	\label{fig:model_fit}
\end{figure}

As shown in Figure~\ref{fig:model_fit}, the constant SIS model ($\sigma_0 \approx 947\text{ km s}^{-1}$) is unable to match the profile: it underestimates the central velocity dispersion by over $100\text{ km s}^{-1}$ and overestimates the dispersion at large radii. This leads to a poor reduced chi-squared of $\chi_\nu^2 = 3.23$ and a large BIC penalty ($\Delta\text{BIC} = 16.42$), clearly ruling out a flat isothermal model for Coma.

The Soft-Core model provides a better fit than the SIS ($\chi_\nu^2 = 1.24$), but it still underpredicts the central dispersion. The NFW profile yields the lowest residuals across the entire range, with $\chi_\nu^2 = 0.69$ and $\text{BIC} = 8.27$. The difference of $\Delta\text{BIC} = 3.36$ shows positive evidence favoring the NFW profile over the Soft-Core profile. These results demonstrate that the observed velocities in Coma favor a cuspy halo profile over a cored one.

\section{Dynamical Mass Estimation and Literature Comparison} \label{sec:mass_estimation}

\subsection{NFW Enclosed Mass and Virial Mass}
Using the best-fit NFW parameters, the enclosed mass profile $M(<r)$ can be calculated by integrating the NFW potential:
\begin{equation}
	M(<r) = M_0 \left[\ln\left(1 + \frac{r}{r_s}\right) - \frac{r/r_s}{1 + r/r_s}\right]
\end{equation}
where the mass normalization $M_0$ is given by:
\begin{equation}
	M_0 = \frac{r_s \cdot \sigma_{\text{max}}^2}{G \cdot f_V}
\end{equation}
Here, $G$ is the gravitational constant and $f_V \approx 0.2162$ is a numerical factor corresponding to the maximum of the velocity dispersion function.

Assuming a cluster redshift of $z = 0.0231$ and a Hubble constant of $H_0 = 70\text{ km s}^{-1}\text{Mpc}^{-1}$ puts Coma at a distance of approximately $D \approx 98.9\text{ Mpc}$. At this distance, the fitted scale angle $\theta_s = 1.047^\circ \pm 0.353^\circ$ corresponds to a physical scale radius of:
\begin{equation}
	r_s \approx 1.81\text{ Mpc} \quad (1810\text{ kpc})
\end{equation}

\begin{figure}[ht!]
	\centering
	\plotone{"Enclosed Mass Profile.png"}
	\caption{Enclosed dynamical mass profile $M(<\theta)$ of the Coma Cluster as a function of projected radius. The dotted vertical line marks the scale radius $r_s \approx 1.05^\circ$ ($1.81\text{ Mpc}$). At the virial radius ($r \approx 2.0\text{ Mpc}$), the enclosed mass is $(5.11 \pm 0.45) \times 10^{14}\,M_\odot$.}
	\label{fig:mass_profile}
\end{figure}

Evaluating $M(<r)$ at the nominal virial radius of Coma ($r_{\text{vir}} \approx 2.0\text{ Mpc}$, corresponding to $R_{200}$) yields a total virial mass of:
\begin{equation}
	M_{\text{vir}} = (5.11 \pm 0.45) \times 10^{14}\,M_\odot
\end{equation}
Within our full $2.0^\circ$ survey aperture ($r \approx 3.45\text{ Mpc}$), the total enclosed mass is $M(<2.0^\circ) = 9.57 \times 10^{14}\,M_\odot$. Figure~\ref{fig:mass_profile} shows the resulting enclosed mass profile as a function of radius.

\subsection{Comparison with Literature}
It is helpful to compare this dynamical mass estimate with previous measurements from optical galaxy kinematics, X-ray gas, and gravitational lensing:

\begin{itemize}
	\item \textbf{Optical Kinematics:} Historically, \citet{kent1982} measured $M_{\text{vir}} \sim (5\text{--}10) \times 10^{14} h^{-1} M_\odot$ using roughly 300 galaxy velocities. More recently, \citet{lokas2003} applied Jeans modeling to 445 galaxies and found $M_{\text{vir}} = (5\text{--}7) \times 10^{14} M_\odot$. Our result of $(5.11 \pm 0.45) \times 10^{14} M_\odot$ is fully consistent with these earlier works, while reducing statistical uncertainties by using a much larger sample of 2,098 galaxies.

	\item \textbf{X-ray Hydrostatic Masses:} X-ray studies using *ROSAT*, *XMM-Newton*, and *Chandra* have measured the temperature and density of the hot intra-cluster gas, yielding hydrostatic masses of $M_{200} \approx (5.0\text{--}6.5) \times 10^{14} M_\odot$ \citep[e.g.,][]{reiprich2002, falco2014}. The close match between our kinematic mass and the X-ray gas mass indicates that Coma is in approximate dynamical equilibrium within $r \le 2.0\text{ Mpc}$.

	\item \textbf{Weak Lensing and Sunyaev-Zel'dovich Effect:} Weak gravitational lensing studies \citep[e.g.,][]{okabe2010} and thermal Sunyaev-Zel'dovich (SZ) measurements from the *Planck* satellite \citep{planck2014} report virial masses in the range of $(5.5\text{--}7.5) \times 10^{14} M_\odot$. Our dynamical result aligns comfortably with these independent mass estimates.
\end{itemize}

\subsection{Discussion of Assumptions and Substructure}
Two main assumptions in this analysis should be noted:
\begin{enumerate}
	\item \textbf{Orbital Isotropy ($\beta = 0$):} In our Jeans analysis, we assumed isotropic galaxy orbits ($\beta(r) = 0$). In real clusters, galaxy orbits may become slightly radial at larger radii, which can steepen the projected velocity dispersion profile. However, because the isotropic NFW model already fits the data well ($\chi_\nu^2 = 0.69$), orbital anisotropy appears to be a modest secondary effect within $r \le 2.0\text{ Mpc}$.
	\item \textbf{Substructure and Infall:} Coma is known to host an infalling sub-cluster centered on NGC~4839, situated roughly $40$ arcminutes southwest of the main cluster core. In our dispersion profile, this appears as the slight increase in dispersion around $\theta \approx 1.3^\circ\text{--}1.6^\circ$ ($910\text{--}930\text{ km s}^{-1}$). Although this subgroup represents an ongoing merger, it is a localized feature and does not disrupt the overall NFW gravitational potential fit.
\end{enumerate}

\section{Conclusions} \label{sec:conclusions}

In this paper, I analyzed spectroscopic data from the DESI Iron data release to examine the kinematics and dark matter halo of the Coma Cluster (Abell 1656). The main conclusions are:

\begin{enumerate}
	\item Using DESI quality flags, relativistic velocity calculations, and $3\sigma$ velocity clipping, I isolated a sample of 2,098 confirmed member galaxies within $\theta \le 2.0^\circ$.
	\item The intrinsic line-of-sight velocity dispersion decreases steadily from $1075.94 \pm 40.72\text{ km s}^{-1}$ in the core to $807.01 \pm 85.07\text{ km s}^{-1}$ at $r \approx 3.2\text{ Mpc}$, with a slight localized increase around $1.4^\circ$ due to the infalling NGC~4839 subgroup.
	\item Comparing three halo models shows that the cuspy NFW profile provides the best fit to the data ($\chi_\nu^2 = 0.69$, $\text{BIC} = 8.27$), while a flat Singular Isothermal Sphere is strongly excluded ($\Delta\text{BIC} = 16.42$) and a Soft-Core profile is disfavored ($\Delta\text{BIC} = 3.36$).
	\item Integrating the NFW profile yields a physical scale radius of $r_s \approx 1.81\text{ Mpc}$ and a virial mass of $M_{\text{vir}} = (5.11 \pm 0.45) \times 10^{14}\,M_\odot$ at $r = 2.0\text{ Mpc}$. This estimate agrees well with independent X-ray and weak lensing measurements.
\end{enumerate}

These results confirm that with large spectroscopic surveys like DESI, cluster galaxy dynamics can place tight constraints on dark matter mass profiles, and that the Coma Cluster continues to be consistent with the cuspy halo predictions of $\Lambda\text{CDM}$ cosmology.

\section*{Acknowledgments}
This work made use of public data from the Dark Energy Spectroscopic Instrument (DESI) Iron data release.

\software{Astropy \citep{astropy2013}, NumPy \citep{harris2020}, SciPy \citep{virtanen2020}, Matplotlib \citep{hunter2007}, pandas.}

\begin{thebibliography}{}
	\bibitem[Astropy Collaboration et al.(2013)]{astropy2013} Astropy Collaboration, et al. 2013, \aap, 558, A33
	\bibitem[DESI Collaboration et al.(2023)]{desi2023} DESI Collaboration, et al. 2023, arXiv:2306.06308
	\bibitem[Falco et al.(2014)]{falco2014} Falco, M., Hansen, S. H., Wojtak, R., et al. 2014, \mnras, 442, 1887
	\bibitem[Harris et al.(2020)]{harris2020} Harris, C. R., et al. 2020, \nat, 585, 357
	\bibitem[Hunter(2007)]{hunter2007} Hunter, J. D. 2007, Computing in Science \& Engineering, 9, 90
	\bibitem[Kent \& Gunn(1982)]{kent1982} Kent, S. M., \& Gunn, J. E. 1982, \aj, 87, 945
	\bibitem[Łokas \& Mamon(2003)]{lokas2003} Łokas, E. L., \& Mamon, G. A. 2003, \mnras, 343, 401
	\bibitem[Navarro et al.(1997)]{navarro1997} Navarro, J. F., Frenk, C. S., \& White, S. D. M. 1997, \apj, 490, 493
	\bibitem[Okabe et al.(2010)]{okabe2010} Okabe, N., Takada, M., Umetsu, K., et al. 2010, \pasj, 62, 811
	\bibitem[Planck Collaboration(2014)]{planck2014} Planck Collaboration, et al. 2014, \aap, 571, A20
	\bibitem[Reiprich \& Böhringer(2002)]{reiprich2002} Reiprich, T. H., \& Böhringer, H. 2002, \apj, 567, 716
	\bibitem[Virtanen et al.(2020)]{virtanen2020} Virtanen, P., et al. 2020, Nature Methods, 17, 261
	\bibitem[Zwicky(1933)]{zwicky1933} Zwicky, F. 1933, Helvetica Physica Acta, 6, 110
\end{thebibliography}

\end{document}
